\documentclass[12pt]{article}
\usepackage[utf8]{inputenc}
\usepackage[T1]{fontenc}
\usepackage{geometry}
\usepackage{hyperref}
\usepackage{enumitem}
\usepackage{tabularx}
\usepackage{colortbl}
\usepackage{xcolor}
\usepackage{booktabs}
\usepackage{titlesec}

% Page setup
\geometry{a4paper, margin=1in}
\hypersetup{
	colorlinks=true,
	linkcolor=blue,
	urlcolor=blue,
	citecolor=blue,
}

% Custom colors
\definecolor{darkblue}{rgb}{0.1,0.2,0.5}
\definecolor{lightgray}{gray}{0.95}

% Section formatting
\titleformat{\section}{\large\bfseries\color{darkblue}}{}{0em}{}
\titleformat{\subsection}{\normalsize\bfseries\color{blue}}{}{0em}{}

\begin{document}
	
	% Title
	\begin{center}
		{\Huge\bfseries Training Program for Big Data \& Distributed Systems} \\[0.5cm]
		{\large Bridge Role — Supporting Until Permanent Position} \\[0.3cm]
		{\textit{Structured, Intensive, Career-Level Training}} \\[1cm]
	\end{center}
	
	\begin{quote}
		\emph{This is a bridge role intended to support you until you secure a permanent position. The training path is structured and suggestive. Realistically, it takes 6 months to 2 years to gain full readiness for the role, depending on}
		
		\vspace{0.3cm}
		
		Company QcFinance.In \\
		Implementation for Joshi LLC NJ USA
		
		\vspace{0.3cm}
		
		\textbf{Additional Learning Resources:} \\
		In addition to the Udemy courses listed below, there are \textbf{YouTube playlists} based on the project that you must follow. These will be shared incrementally as you progress through each phase.
		
		\vspace{0.3cm}
		
		The courses must be taken in the order listed below, as each builds on the previous one. Training will focus heavily on \textbf{Python}, \textbf{distributed systems}, \textbf{testing}, and \textbf{big-data computation}.
		
		Let me know once you start, and we will review progress weekly.
		
		Your role might not be exactly what is in the course, but it will give you some idea of what skills are needed and how to approach them. As time goes by, the role might evolve and become more focused on Copilot and AI-based prompts and agents. However, the working style, methodology, and culture will remain the same.
		
		Best regards, \\
		Shivgan Joshi India MP Indore
	\end{quote}
	
	\newpage
	
	% ====================================================
	% COURSE LINKS (COMPLETE LIST — ALL 13 COURSES)
	% ====================================================
	\section*{Complete Course Links (Suggested Sequence)}
	
	\begin{enumerate}[label=\textbf{Course \arabic*:}]
		
		\item Note taking Course and Shadow Analyst for C Suite Analyst\\
		\href{https://www.udemy.com/course/note-taking-course-and-shadow-analyst-for-c-suite-analyst/}
		{https://www.udemy.com/course/note-taking-course-and-shadow-analyst-for-c-suite-analyst/}
		
		\item Remote \& Independent Working in Managerless Environments \\
		\href{https://www.udemy.com/course/remote-independent-working-in-a-managerless-environment-101/}
		{https://www.udemy.com/course/remote-independent-working-in-a-managerless-environment-101/}
		
		\item \textbf{Connecting \& Managing Ubuntu Bridge Computers for Remote Work} \\
		\href{https://www.udemy.com/course/connecting-managing-ubuntu-bridge-computers-for-remote-work/}
		{https://www.udemy.com/course/connecting-managing-ubuntu-bridge-computers-for-remote-work/}
		
		\item Big Data Infrastructure Administration — Ticket-Based Learning \\
		\href{https://www.udemy.com/course/big-data-infra-admin-101-ticket-based-learning/}
		{https://www.udemy.com/course/big-data-infra-admin-101-ticket-based-learning/}
		
		\item Python Full-Stack Computational Framework (Beginner) \\
		\href{https://www.udemy.com/course/python-full-stack-computational-framework-for-begineers-101/}
		{https://www.udemy.com/course/python-full-stack-computational-framework-for-begineers-101/}
		
		\item Full-Spectrum Comprehensive Testing \& Error Reading \\
		\href{https://www.udemy.com/course/full-spectrum-comprehensive-testing-error-reading-101/}
		{https://www.udemy.com/course/full-spectrum-comprehensive-testing-error-reading-101/}
		
		\item Research \& Writing Skills for Essays in Controlled Environment on Bridge Node\\
		\href{https://www.udemy.com/course/research-writing-skills-for-essays-in-controlled-environment}
		{https://www.udemy.com/course/research-writing-skills-for-essays-in-controlled-environment}
		
		\item Python Big Data Refresher — Intermediate (Quiz Based) \\
		\href{https://www.udemy.com/course/refresher-python-big-data-intermediate-quiz-based-course-102/}
		{https://www.udemy.com/course/refresher-python-big-data-intermediate-quiz-based-course-102/}
		
		\item Python Full-Stack Backend \& ML Engines \\
		\href{https://www.udemy.com/course/python-full-stack-and-backend-engines-for-mc-ml-engines-102/}
		{https://www.udemy.com/course/python-full-stack-and-backend-engines-for-mc-ml-engines-102/}
		
		\item Writing Tests for Simulation Engineering \& Code Conversion \\
		\href{https://www.udemy.com/course/writing-tests-for-simeng-python-code-conversion-concepts-101/}
		{https://www.udemy.com/course/writing-tests-for-simeng-python-code-conversion-concepts-101/}
		
		\item Full-Stack Computational Framework for Big Data Simulations \\
		\href{https://www.udemy.com/course/full-stack-computational-framework-for-big-data-simulations/}
		{https://www.udemy.com/course/full-stack-computational-framework-for-big-data-simulations/}
		
		\item \textbf{Big Data Code Changes for Full-Stack Simulation Engine 2026} \\
		\href{https://www.udemy.com/course/big-data-code-changes-for-full-stack-simulation-engine-2026/}
		{https://www.udemy.com/course/big-data-code-changes-for-full-stack-simulation-engine-2026/}
		
		\item Converting Standalone Python Code to Distributed Analytics Code \\
		\href{https://www.udemy.com/course/converting-standalone-code-to-distributed-code-for-analytics/}
		{https://www.udemy.com/course/converting-standalone-code-to-distributed-code-for-analytics/}
	\end{enumerate}
	
	% ====================================================
	% ASSUMPTIONS
	% ====================================================
	\section*{Training Assumptions}
	
	\begin{itemize}
		\item Each course requires approximately 2–4 days
		\item Average effort: 4–6 hours per day
		\item Training cadence: one course per week
		\item Initial readiness: approximately 13 weeks (all 13 courses)
		\item Professional mastery requires extended real-world practice
	\end{itemize}
	
	\newpage
	
	% ====================================================
	% DETAILED TRAINING PLAN (ALL 13 COURSES)
	% ====================================================
	\section*{Part 2: Detailed Training Plan (Step-by-Step, Time-Based)}
	
	\subsection*{Assumptions}
	\begin{itemize}
		\item Each course takes 2–4 days
		\item Average effort: 4–6 hours/day
		\item Training pace: 1 course per week
		\item Total timeline below focuses on initial readiness (approx. 13 weeks)
		\item Full mastery continues beyond this with practice
	\end{itemize}
	
	\vspace{0.5cm}
	
	% PHASE 1 — Foundation, Shadow Analyst & Bridge Setup
	\section*{Phase 1: Foundation, Shadow Analyst \& Bridge Setup (Weeks 1–3)}
	
	\begin{tabularx}{\textwidth}{|X|p{3cm}|p{3cm}|X|}
		\hline
		\textbf{Week} & \textbf{Course} & \textbf{Time} & \textbf{Goal} \\
		\hline
		1 & \href{https://www.udemy.com/course/note-taking-course-and-shadow-analyst-for-c-suite-analyst/}{Course 1: Note taking \& Shadow Analyst for C Suite} & 2–3 days & Learn executive-level note taking, shadowing techniques, and how to support C-suite decision-making with structured observations \\
		\hline
		2 & \href{https://www.udemy.com/course/remote-independent-working-in-a-managerless-environment-101/}{Course 2: Remote \& Independent Working} & 2–3 days & Learn how to work independently; understand ticket‑based and manager‑less environments; set expectations for self‑driven technical work \\
		\hline
		3 & \href{https://www.udemy.com/course/connecting-managing-ubuntu-bridge-computers-for-remote-work/}{Course 3: Ubuntu Bridge Computers for Remote Work} & 3–4 days & Set up and manage Ubuntu-based bridge machines; understand networking, SSH, and remote connectivity essential for distributed infrastructure \\
		\hline
	\end{tabularx}
	
	\vspace{0.3cm}
	
	% PHASE 2 — Infrastructure & Core Python
	\section*{Phase 2: Infrastructure \& Core Python (Weeks 4–5)}
	
	\begin{tabularx}{\textwidth}{|X|p{3cm}|p{3cm}|X|}
		\hline
		\textbf{Week} & \textbf{Course} & \textbf{Time} & \textbf{Goal} \\
		\hline
		4 & \href{https://www.udemy.com/course/big-data-infra-admin-101-ticket-based-learning/}{Course 4: Big Data Infrastructure \& Ticket‑Based Learning} & 3–4 days & Understand big‑data systems; learn production‑style workflows; exposure to infrastructure issues and troubleshooting \\
		\hline
		5 & \href{https://www.udemy.com/course/python-full-stack-computational-framework-for-begineers-101/}{Course 5: Python Full‑Stack Computational Framework (Beginner)} & 3–4 days & Python fundamentals; data structures and computational thinking; writing clean, reusable Python code \\
		\hline
	\end{tabularx}
	
	\vspace{0.3cm}
	
	% PHASE 3 — Testing, Research, Intermediate Python
	\section*{Phase 3: Testing, Research Writing \& Intermediate Python (Weeks 6–9)}
	
	\begin{tabularx}{\textwidth}{|X|p{3cm}|p{3cm}|X|}
		\hline
		\textbf{Week} & \textbf{Course} & \textbf{Time} & \textbf{Goal} \\
		\hline
		6 & \href{https://www.udemy.com/course/full-spectrum-comprehensive-testing-error-reading-101/}{Course 6: Comprehensive Testing \& Error Reading} & 2–3 days & Debugging skills; understanding stack traces; reading and fixing errors independently \\
		\hline
		7 & \href{https://www.udemy.com/course/research-writing-skills-for-essays-in-controlled-environment}{Course 7: Research \& Writing in Controlled Environment} & 2–3 days & Structured documentation, clear technical writing, and research skills for controlled big-data environments \\
		\hline
		8 & \href{https://www.udemy.com/course/refresher-python-big-data-intermediate-quiz-based-course-102/}{Course 8: Python Big Data Refresher (Intermediate)} & 3–4 days & Reinforce Python concepts; apply Python to data‑heavy problems; build confidence in writing intermediate‑level code \\
		\hline
		9 & \href{https://www.udemy.com/course/python-full-stack-and-backend-engines-for-mc-ml-engines-102/}{Course 9: Python Backend \& ML Engines} & 3–4 days & Backend computation; understanding Python in ML / analytics pipelines; performance‑oriented programming \\
		\hline
	\end{tabularx}
	
	\vspace{0.3cm}
	
	% PHASE 4 — Testing at Scale, Simulation Code & Distributed Systems
	\section*{Phase 4: Testing at Scale, Simulation Code \& Distributed Systems (Weeks 10–13)}
	
	\begin{tabularx}{\textwidth}{|X|p{3cm}|p{3cm}|X|}
		\hline
		\textbf{Week} & \textbf{Course} & \textbf{Time} & \textbf{Goal} \\
		\hline
		10 & \href{https://www.udemy.com/course/writing-tests-for-simeng-python-code-conversion-concepts-101/}{Course 10: Writing Tests for Simulation \& Code Conversion} & 2–3 days & Writing tests for complex logic; ensuring correctness during refactoring; quality control mindset \\
		\hline
		11 & \href{https://www.udemy.com/course/full-stack-computational-framework-for-big-data-simulations/}{Course 11: Full‑Stack Computational Framework for Big Data} & 3–4 days & End‑to‑end computation pipelines; simulation‑based design; scaling logic from small to large systems \\
		\hline
		12 & \href{https://www.udemy.com/course/big-data-code-changes-for-full-stack-simulation-engine-2026/}{Course 12: Big Data Code Changes for Full-Stack Simulation Engine 2026} & 3–4 days & Hands-on modifying simulation engine code for big-data scale; real-world code change management \\
		\hline
		13 & \href{https://www.udemy.com/course/converting-standalone-code-to-distributed-code-for-analytics/}{Course 13: Converting Standalone Code to Distributed Code} & 3–4 days & Distributed computing concepts; parallel execution; production‑grade analytics workflows \\
		\hline
	\end{tabularx}
	
	\newpage
	
	% SUMMARY TIMELINE — UPDATED (13 COURSES)
	\section*{Summary Timeline (Updated)}
	
	\begin{tabularx}{\textwidth}{|l|X|}
		\hline
		\textbf{Phase} & \textbf{Duration} \\
		\hline
		Foundation, Shadow Analyst, Remote Work \& Bridge Setup & 3 weeks \\
		\hline
		Core Python \& Infrastructure & 2 weeks \\
		\hline
		Testing, Research Writing \& Intermediate Python & 4 weeks \\
		\hline
		Distributed Systems, Simulation Code \& Scale & 4 weeks \\
		\hline
		\textbf{Initial Readiness} & \textbf{\~13 weeks} \\
		\hline
		Professional Mastery & 6–24 months with practice \\
		\hline
	\end{tabularx}
	
	\vspace{0.5cm}
	
	\section*{How to Use This Plan}
	\begin{itemize}
		\item Follow strict order
		\item Do hands‑on coding alongside videos
		\item Re‑do exercises if unclear
		\item Expect difficulty — that’s normal
		\item This is career‑level training, not casual learning
	\end{itemize}
	
	\newpage
	
	% ====================================================
	% EXPANDED ACRONYMS AND PEDAGOGICAL BASIS — UPDATED
	% ====================================================
	\section*{Appendix: Expanded Acronyms \& Pedagogical Basis}
	
	\subsection*{Expanded Acronyms (Full Forms)}
	\begin{itemize}
		\item \textbf{ML} — Machine Learning
		\item \textbf{MC} — Monte Carlo (simulation methods)
		\item \textbf{SimEng} — Simulation Engineering
		\item \textbf{Infra} — Infrastructure
		\item \textbf{Admin} — Administration
		\item \textbf{Python} — Not an acronym; a programming language named after Monty Python
		\item \textbf{API} — Application Programming Interface
		\item \textbf{CI/CD} — Continuous Integration / Continuous Deployment
		\item \textbf{SSH} — Secure Shell
		\item \textbf{Ubuntu Bridge} — A computer acting as a network bridge/gateway for remote distributed systems
	\end{itemize}
	
	\subsection*{Why This Strict Order? (Pedagogical Explanation — Updated for 13 Courses)}
	
	The training plan follows a \textbf{cumulative skill ladder}.
	
	\begin{enumerate}
		\item \textbf{Week 1 (Shadow Analyst)} — Executive communication and observation skills; foundation for understanding business context.
		
		\item \textbf{Week 2 (Remote Work)} — Metacognitive foundation: how to learn independently.
		
		\item \textbf{Week 3 (Ubuntu Bridge Computers)} — Before any cloud or big-data infra, you must know how to set up and manage the actual bridge machine (Ubuntu) that connects remote work environments. Essential for the "bridge role."
		
		\item \textbf{Week 4 (Big Data Infrastructure)} — Understand ecosystems (clusters, pipelines) before writing code.
		
		\item \textbf{Week 5 (Python Beginner)} — Core programming fundamentals.
		
		\item \textbf{Week 6 (Testing \& Error Reading)} — Debugging taught early for self-correction.
		
		\item \textbf{Week 7 (Research \& Writing)} — Structured documentation and controlled-environment research skills. Critical for tickets and simulation reports.
		
		\item \textbf{Week 8 (Python Intermediate Refresher)} — Reinforcement through quiz-based, data-heavy problems.
		
		\item \textbf{Week 9 (Backend \& ML Engines)} — Performance-oriented Python for real pipelines.
		
		\item \textbf{Week 10 (Writing Tests)} — Non-negotiable for distributed systems.
		
		\item \textbf{Week 11 (Full-Stack Big Data Framework)} — Integrates all previous skills.
		
		\item \textbf{Week 12 (Big Data Code Changes for Simulation Engine)} — Real-world modification of existing simulation engine code. Teaches how to safely introduce big-data changes without breaking simulations.
		
		\item \textbf{Week 13 (Standalone → Distributed)} — Capstone: convert single-machine code to parallel, distributed execution.
	\end{enumerate}
	
	\subsection*{Cognitive Load Management}
	\begin{itemize}
		\item Early weeks focus on \emph{one new major concept} per week.
		\item Shadow analyst skills establish business communication before technical depth.
		\item Ubuntu bridge skills are introduced immediately after remote work — practical before theoretical infrastructure.
		\item Research/writing is placed after testing to reinforce documentation of errors and test results.
		\item Big-data code changes come before final conversion to distributed code — you first learn to modify existing engines, then redesign from standalone to distributed.
	\end{itemize}
	
	\subsection*{Time-to-Competence Realism (Updated for 13 Courses)}
	\begin{itemize}
		\item \textbf{13 weeks} → Initial readiness (can complete supervised tickets, understands bridge role basics)
		\item \textbf{6 months} → Consistent independent work on medium‑complexity tasks
		\item \textbf{1–2 years} → Full mastery (architecting distributed solutions, mentoring others)
	\end{itemize}
	
	\vspace{1cm}
	\noindent\textit{Document generated from the original training plan by Shivgan Joshi. All Udemy links are preserved as provided. Updated with all 13 courses in strict sequential order.}
	
\end{document}